\section{从保留记录到分析单位}
\label{sec:materials}
\subsection{任务、候选、审核记录与证据签名}
\label{subsec:record-units}
\emph{任务}是一项数学作业，例如形式化某个指定定义，或证明某条指定定理。\emph{候选}是为该任务提交并保存的一个版本。同一任务可以有多个候选，同一候选也可以被审核多次。\emph{审核记录}保留对一个对象的判断，以及能够取得的相关证据和上下文；这个对象可能是候选，也可能是正式成果。

分析使用\emph{证据签名}标识保留下来的、由审核者产生的证据。整理后的主要集合包含 3,231 个这样的签名。这是按档案识别规则计算的证据条目数，不是题目数，也不保证发生了 3,231 次彼此分开的审核执行。例如，同一次审核的原始结果文件和整理结果文件，可能作为不同证据条目保留下来。后续分类会识别其中一些关系，而不把它们解释为独立重复审核。

\subsection{一道题、三次提交、两条连接、一个片段}
\label{subsec:running-example}
考虑一个说明性定理任务，先后有甲、乙、丙三个提交版本，相应审核记录为 $r_A,r_B,r_C$，判断依次是失败、失败、通过。假设甲到乙、乙到丙都符合下面将定义的候选修改规则，而且两条连接的中间端点确实是同一条审核记录 $r_B$。这个任务就提供了：
\begin{itemize}[leftmargin=*,nosep]
\item 一道题和三个送审候选版本；
\item 本例中的三条审核记录；
\item 两条有向\emph{连接}，也称\emph{边}：$r_A\to r_B$ 和 $r_B\to r_C$；
\item 在 $r_B$ 处拼接这两条边得到的一个\emph{连续片段}。
\end{itemize}
因此，两条边共有四次端点出现，却只有三条不同的端点记录。重复使用中间记录，既影响重建，也影响统计依赖。连接描述两个保留观察之间的关系，并不计算其间发生了多少次编辑或模型调用。

连续性要求共享同一条审核记录，仅有相同题目标识或相同候选文本还不够。假设后来的丁版本之前出现了正式成果复查、不符合选定规则的连接，或时间顺序尚未解决的缺口，就不能仅因题目未变而把丁接入原片段。如果丁与更晚的戊版本又形成一段合格序列，同一道题便贡献第二个片段。这种构造没有断言实际工作曾在缺口处停止。

完整片段与针对某个问题实际观察的部分也有区别。在“失败、失败、通过”的例子中，观察到首次通过会使用两条边；在另一个说明性的“失败、通过、失败”片段中，完整重建仍可能包含两条合格边，但首次通过分析只使用第一条。它不会把最后的失败当成通过之前的失败。这些例子用于解释规则，不是新增观察。

\subsection{从证据到 2,645 条可定向相邻连接}
\label{subsec:adjacency}
为了研究审核之间发生了什么变化，分析从档案已经固定的同任务相邻关系出发。只有一对原始相邻记录的两端签名都属于整理后的主要集合时，才保留这对记录。这里\emph{不会}先删掉中间记录，再把剩余记录中变成邻居的两条重新连接起来；原有缺口因此仍然保留。

这项成员规则得到 2,738 对记录。其中 93 对按已声明的时间戳规则无法确定先后，包括时间相同或精度不足等情况；它们仍作为无向证据保留，不进入前后比较。其余 $2,738-93=2,645$ 对提供有向连接。这个方向是根据记录时间和标识建立的顺序代理，不证明一次审核导致了下一事件。也不能简单用 3,231 减去题目数得到连接数，因为规则保留的是原始相邻关系，而不是筛选后重新搭成无缺口链条。

即使连接已有方向，也不一定表示修改。前一端可能审核候选，后一端审核文本相同的正式成果；两端也可能都审核正式成果；原始结果文件之后还可能只是同次审核的整理版本。因此，下一步需要根据审核对象、主要 Lean 文本、来源、已登记上下文，以及此前的结构化问题来分类活动。通过或失败标签本身不用于决定这个活动类别。

\subsection{哪些连接算候选修改：465 条与 367 条}
\label{subsec:revision-rule}
分类程序按顺序应用规则：先识别同次审核的两种文件表示，再识别明确标记的外部审核，随后检查相同登记条件下是否有两次独立执行的明确证据。条件相同却无法证明执行分开的情况仍归为未解决。程序接着识别候选以相同主要文本落地到正式成果的情况。处理过这些优先类别后，才判定\emph{候选修改}：后一审核对象必须是候选，而且两端的主要 Lean 文本不同。剩余的正式成果到正式成果连接归为正式复查，其他无法解决的连接保留为混合或不明类别。

在候选修改类别内，\emph{高置信}表示前一审核包含结构化问题；没有这种问题记录时，标为\emph{中置信}。这是可以重复计算的规则标签，不是数值置信水平，也不认证某个数学问题已经修复。规则识别的是记录问题之后出现了改变的候选，并未独立确认改变解决了那些问题。程序将这一类别命名为 \code{candidate\_repair}；正文使用“候选修改”，以保留这个区别。

最终类别对全部有向连接形成如下划分：
\begin{table}[htbp]
\centering\small
\caption{2,645 条有向连接的活动分类。计数单位是连接，不是执行次数或题目。}
\label{tab:roles}
\begin{tabularx}{\linewidth}{Y r Y}
\toprule
类别 & 数量 & 含义 \\
\midrule
候选修改 & 465 & 主要文本改变，后一对象为候选 \\
候选落地 & 241 & 候选之后是主要文本相同的正式成果 \\
正式成果复查 & 1,176 & 两端对象均为正式成果 \\
外部审核 & 207 & 至少一端来源明确标记为外部审核 \\
同次审核的两种表示 & 166 & 已识别的原始与整理结果配对 \\
同条件的不同审核 & 0 & 需要分别执行的明确证据 \\
混合或不明 & 390 & 现有字段不足以确定活动 \\
\midrule
合计 & 2,645 & 每条连接只归入一类 \\
\bottomrule
\end{tabularx}
\end{table}

465 条候选修改边涉及 250 个不同任务，分为 367 条高置信边和 98 条中置信边。完整的 465 条边用于第~\ref{sec:statistics}~节的修改配对分析。其中，高置信敏感性分析使用 367 条边，涉及 207 个任务。任务数统计的是筛选后出现的不同标识；与互不重叠的边数不同，各子集的任务数未必可以相加，因为同一个任务可能同时贡献两类边。

\subsection{把 367 条边拼成 262 个完整连续片段}
\label{subsec:segment-rule}
过程问题的重建只使用 367 条高置信候选修改边。从没有合格前驱的边出发，仅当前一条边的较后审核标识，与下一条边的较早审核标识相同时，才追加那条合格边，直到找不到下一条。结果称为\emph{极大连续片段}，即在所选边集合内已经无法继续延伸的链。档案称它为片段，后面的章节也使用\emph{路径}或\emph{过程}。

这样的归组让每条选定边恰好进入一个片段，共得到 262 个片段。遇到不合格连接、时间顺序未解决的边界，或已观察连接结束时，片段便结束。记录通过本身不终止这次重建：通过标签将在下一项分析操作中使用。表~\ref{tab:segment-lengths}~展示边与片段之间的数量对应。

\begin{table}[htbp]
\centering\small
\caption{完整重建片段的长度。此时尚未筛选失败起点，也尚未在首次通过处停止。}
\label{tab:segment-lengths}
\begin{tabular}{rrr}
\toprule
每片段边数 & 片段数 & 贡献的边数 \\
\midrule
1 & 212 & 212 \\
2 & 38 & 76 \\
3 & 1 & 3 \\
4 & 5 & 20 \\
5 & 3 & 15 \\
12 & 2 & 24 \\
17 & 1 & 17 \\
\midrule
合计 & 262 & 367 \\
\bottomrule
\end{tabular}
\end{table}

第二列数片段，第三列数片段所含的边。归组本身没有删除任何边，所以 262 个片段仍涉及原来 367 条边对应的 207 个任务。同任务的多个片段在重建中仍是分开的路径，在统计归组时则仍由同一任务标识联系起来。

\subsection{选出 230 个失败起点，再观察 329 个步骤}
\label{subsec:failure-cohort}
过程分析询问明确记录失败之后发生了什么，所以选取第一条审核明确写为失败的 230 个片段，把另外 32 个排除在这个问题之外。选中的片段属于 182 个不同任务。230 数的是起始片段，182 数的是它们涉及的任务标识；两者都不是独立审核执行的数量。

现在，在每个选中片段内，观察到第一个较后审核记录通过时停止；如果没有出现通过，就在片段边界停止。这第二项操作改变的是已观察长度，不改变此前固定的片段成员。对选中片段应用该规则，共得到 329 个观察步骤。每一步是一条合格边，结果取自它的较后审核。230 条路径中，206 条以首次记录通过结束，24 条到片段边界时仍未见通过。原始失败记录是起始条件，不额外算作一个观察步骤。

因此，329 不是把全部 367 条重建边换一个名字：其间经过了失败起点筛选与首次通过截取。同样，$206/230$ 不是 182 个任务中的成功比例，因为同一任务可以贡献多条路径。第~\ref{sec:process}~节将继续解释逐步变化的分母，并区分路径退出与任务永久失败。

\begin{table}[htbp]
\centering\small
\caption{活动分类之后的两条历史分析路线。每行均说明操作及其计数单位。}
\label{tab:units}
\begin{tabularx}{\linewidth}{Y Y}
\toprule
操作 & 结果与用途 \\
\midrule
保留全部候选修改连接 & 465 条边、250 个任务：配对的记录标签变化 \\
进一步限于高置信边 & 367 条边、207 个任务：配对敏感性分析，以及路径重建的输入 \\
在共享审核端点处拼接高置信边 & 262 个完整片段：先确定连续性，再按结果筛选 \\
保留明确失败起点的片段 & 230 个片段、182 个任务：固定的过程分析集合 \\
沿每个选中片段观察到首次通过或边界 & 329 个步骤；206 个通过终点和 24 次退出 \\
\bottomrule
\end{tabularx}
\end{table}

\subsection{为什么选定执行样本是 14 对候选}
\label{subsec:execution-sample}
较大历史描述记录判断。为了考察明确规定的评价器能够确认代码的哪些性质，研究另外从修改材料中挑选一份较小的候选配对样本：6 对开发候选与 8 对按任务分开的留出候选，共 14 对、28 个候选端点。同一任务不会拆到开发和留出两组。这是回顾性研究中的选取分区；此前已经接触过部分档案，因此不能把它解释为完全盲态、独立抽样的实验。

每对候选接受两种评估。一种检查允许的公理依赖，另一种检查主定理不能直接假设完整结论这一狭窄要求。每项评估都把两个候选分别放在两条明确标准下考察，四种候选与标准的组合构成一张\emph{四格表}，第~\ref{subsec:four-cells}~节会给出定义和完整计算。因此共有 $14\times2=28$ 张历史表，它们重复使用同一批 14 对候选。第~\ref{sec:audit}~节报告测量结果；本版没有增加新的执行。

\subsection{构造比较、定义见证与布尔任务}
\label{subsec:constructed-materials}
研究还询问：在人为安排的已知或缺失分数组合中，比较方法是否按预期工作？7 张构造表服务于这个问题，与 28 张历史表合起来，共有 35 张用于考察比较程序的表。这个总数数的是表，不是 35 个独立抽样任务。第~\ref{sec:audit}~节会解释每张表所提出的六个诊断问题。

另一项历史定义见证考察一道题的两版候选，它们对“可数”采用了不同约定。6 条已证明命题确认候选、明确参照与空集之间的关系；它们描述一个案例，并不表示六次独立成功。

最后，一个有限布尔任务要求当且仅当两个输入不同时输出真。全部输入只有四种，因而可以给出完整行为参照。研究有意设计了 9 种实现：3 种符合要求、5 种不符合要求，以及 1 种编译失败对照。第~\ref{sec:measurement-classes}~节使用其中 8 种可执行实现推导测量的局限，第~\ref{subsec:boolean-audit}~节列出完整结果。一个任务的九种实现，并不提供九个独立抽样数学任务。

这些材料用途相连，却各有自己的分母。历史修改与路径分析描述经过选择的记录；选定执行与构造案例检查具体评价方法，或给出反例。不能把它们的计数相加得到统一样本量，也不能视为同一项统计结论的独立重复验证。
